\section{Metodologia}
\label{sec:metodologia}

A metodologia descreve o ``como'' a pesquisa foi conduzida. Nela, você deve detalhar a arquitetura da solução proposta, os algoritmos criados, os materiais utilizados, os conjuntos de dados (datasets) selecionados e o desenho dos experimentos. O nível de detalhes deve ser suficiente para que outro pesquisador consiga reproduzir o seu trabalho.

Para organizar os dados, resultados parciais ou conceitos, você pode fazer uso de Tabelas, Quadros e Algoritmos. Geralmente, as tabelas apresentam dados numéricos quantitativos (por convenção, com bordas laterais abertas), enquanto os quadros organizam informações textuais e qualitativas (com todas as bordas fechadas). 

A seguir, a Tabela~\ref{tab:exemplo_tabela} apresenta os dados coletados nos experimentos de tempo de execução, inserida através do arquivo contido na pasta \textit{tabelas/}.

\begin{table}[H]
  \centering
  \caption{Exemplo de Tabela (foco em dados numéricos, bordas laterais abertas segundo a ABNT).}
  \label{tab:exemplo_tabela}
  \begin{tabular}{lcc}
    \hline
    \textbf{Métrica} & \textbf{Cenário 1} & \textbf{Cenário 2} \\
    \hline
    Tempo de Execução (ms) & 12 & 15 \\
    Precisão Global (\%) & 98 & 95 \\
    Consumo de Memória (MB) & 100 & 120 \\
    \hline
  \end{tabular}
\end{table}


Já o Quadro~\ref{qua:exemplo_quadro} descreve as ferramentas utilizadas em cada fase do projeto para fins de mapeamento teórico.

\begin{quadro}[H] %htbp
  \centering
  \caption{Exemplo de Quadro (foco em textos ou conceitos, todas as bordas fechadas).}
  \label{qua:exemplo_quadro}
  \resizebox{\textwidth}{!}{
  \begin{tabular}{|l|l|l|}
    \hline
    \textbf{Fase do Projeto} & \textbf{Ferramenta Utilizada} & \textbf{Descrição/Objetivo} \\
    \hline
    Coleta de Dados & Scrapy & Extração automatizada de textos da web \\
    \hline
    Pré-processamento & NLTK (Python) & Remoção de stopwords e tokenização \\
    \hline
    Treinamento & Scikit-Learn & Treinamento do modelo de classificação \\
    \hline
  \end{tabular}
  }
\end{quadro}


Se apropriado, utilize diagramas de fluxo, arquitetura ou pseudocódigos para ilustrar a lógica do seu método, como mostrado no Algoritmo~\ref{alg:exemplo_busca}, que é importado da pasta \textit{algoritmos/}.

\begin{algorithm}[H]%htbp
  \caption{Exemplo de Algoritmo de Busca}
  \label{alg:exemplo_busca}
  \begin{algorithmic}[1]
    \Require Grafo $G$, Nó inicial $s$, Nó objetivo $t$
    \Ensure Caminho percorrido de $s$ até $t$
    \State $Atual \gets s$
    \While{$Atual \neq t$}
      %\State $Vizinhos \gets \text{vizinhos de } Atual$
      \State $Atual \gets \arg\min_{v \in Vizinhos} h(v)$ \Comment{Escolhe o nó mais promissor}
    \EndWhile
    \State \Return Caminho percorrido
  \end{algorithmic}
\end{algorithm}


Caso o seu trabalho seja focado no desenvolvimento de um sistema (como um TCC, por exemplo: \cite{tcc_exemplo}), detalhe aqui as etapas de engenharia de software, modelagem e implementação.
